\capredactadotrue
\section{Plan de Marketing}

El presente plan de marketing establece la estrategia comercial, de comunicación y de penetración de mercado de \nombreproyecto{}. El modelo de negocio se apoya en una arquitectura de red centrada en el \textit{Dungeon Master} (DM): cada DM que adopta la herramienta atrae e integra de forma directa a un grupo promedio de entre cuatro y seis jugadores por mesa. A partir de este principio de difusión, el plan traza un camino de expansión gradual en un horizonte de tres años: comienza con la validación del producto mínimo viable (MVP) en el nodo local de Rosario y escala hacia los principales centros urbanos de Argentina. La consolidación del posicionamiento en el mercado regional del MERCOSUR se considera una línea de evolución posterior a este horizonte.

\subsection{Objetivos}
El objetivo cualitativo del plan es darse a conocer, en un tiempo determinado, entre distintos grupos de jugadores de la región mediante diversos medios, y sentar las bases de una comunidad creciente que genere contenido.

En términos cuantitativos, el plan persigue las metas de cobertura y participación por etapa que se detallan en la sección siguiente: presencia en el 30\,\% de las mesas de rol activas de Rosario durante la etapa de lanzamiento, en el 15\,\% de las mesas del corredor regional durante la de consolidación y en al menos el 20\,\% de las mesas del país durante la de expansión.

\subsection{Resultados esperados en cobertura y participación}

\subsubsection{Cobertura}
La cobertura avanza en tres fases geográficas claramente delimitadas, que coinciden con las etapas del plan.
\begin{itemize}
  \item \textbf{Fase local (etapa de lanzamiento, primer año).} La cobertura inicial se concentra en Rosario y el Gran Rosario. Se utilizan como canales directos de difusión los clubes de juegos de mesa, las agrupaciones universitarias y convenciones como Rosario Juega Rol, Crack Bang Boom y otros eventos menores. La meta es alcanzar presencia en el 30\,\% de las mesas de rol activas de la ciudad, validar la usabilidad y recopilar retroalimentación directa de los primeros DMs adoptantes.
  \item \textbf{Fase provincial y regional (etapa de consolidación, segundo año).} La cobertura se extiende al resto de la provincia de Santa Fe, con foco en la capital provincial y Rafaela, y a las provincias vecinas de la Región Centro y Litoral (Córdoba y Entre Ríos). Esta fase conecta las comunidades de rol mediante eventos de otras provincias y alianzas con tiendas especializadas de la región, con el objetivo de lograr presencia en el 15\,\% de las mesas de ese corredor.
  \item \textbf{Fase nacional (etapa de expansión, tercer año).} Se alcanza la cobertura de todo el territorio argentino, integrando los principales polos urbanos y de mayor densidad de jugadores (CABA, Gran Buenos Aires, Mendoza, Tucumán y Patagonia). La plataforma busca posicionarse en al menos el 20\,\% de las mesas de rol activas del país mediante campañas digitales focalizadas y patrocinios en transmisiones en vivo (\textit{actual plays}) argentinas.
\end{itemize}

\subsubsection{Participación}
\begin{itemize}
  \item \textbf{Etapa de lanzamiento (local):} capturar el 15\,\% del mercado de DMs digitales de la ciudad de Rosario y su área de influencia.
  \item \textbf{Etapa de consolidación (provincial y regional):} alcanzar una cuota del 8\,\% de los DMs activos del corredor Santa Fe, Córdoba y Entre Ríos.
  \item \textbf{Etapa de expansión (nacional):} capturar el 5\,\% del mercado total de DMs digitales de la República Argentina.
\end{itemize}

\subsection{Metas de posicionamiento}
El posicionamiento busca construir una identidad de marca clara y diferenciada en la mente de los DMs de la región, distante de las soluciones anglosajonas tradicionales, mediante atributos de usabilidad, relevancia económica y conexión comunitaria.

\subsubsection{Metas de percepción en el usuario}
\begin{itemize}
  \item \textbf{Eficiencia y reducción de la carga mental.} La meta principal es que el 85\,\% de los DMs encuestados perciba a la aplicación como la solución que más tiempo les ahorra en la preparación previa a la sesión (\textit{pre-session prep}). La plataforma no busca reemplazar la imaginación del director de juego, sino automatizar las tareas administrativas y mecánicas (seguimiento de iniciativa, estadísticas de monstruos, inventario del grupo y notas de mundo) para que pueda concentrarse en la narrativa.
  \item \textbf{Comunidad colaborativa y economía de creadores (\textit{marketplace}).} Se trabaja para posicionar a la plataforma como el principal espacio de intercambio, monetización y colaboración de la comunidad de rol de la región. A través de un \textit{marketplace} integrado, que se habilita en la tercera etapa del plan, la plataforma busca ser percibida bajo un doble estándar de valor:
  \begin{itemize}
    \item \emph{Impulso al talento local:} ser reconocida por los creadores de contenido (autores de aventuras, diseñadores de mapas, ilustradores) como una vitrina justa, que redistribuye directamente a los creadores un porcentaje transparente y competitivo de las ventas y fomenta el desarrollo de la industria TTRPG en América Latina.
    \item \emph{Estándar de contenido propio (Línea Esencial):} ser percibida como un proveedor confiable de módulos, mapas y recursos de uso común y alta calidad, desarrollados por la propia plataforma, que garantizan que cualquier DM, sin importar su presupuesto o su tiempo disponible, cuente con recursos fundamentales listos para usar en sus partidas.
  \end{itemize}
  \item \textbf{Versatilidad de uso (híbrida).} La plataforma busca ser percibida no como un competidor rígido de los tableros virtuales (VTT como Foundry o Roll20), sino como el asistente de gestión definitivo, que funciona de manera transparente tanto para quienes juegan alrededor de una mesa física con papel y lápiz como para quienes utilizan herramientas de transmisión o juego a distancia.
\end{itemize}

\subsection{Producto: estrategia de producto}

\subsubsection{Descripción del producto o servicio (qué es y qué no es)}
\nombreproyecto{} es una herramienta para la organización de la información de campañas, personajes y sesiones de juegos de rol de mesa, disponible como aplicación web para DMs y como aplicación móvil para jugadores. Además de ordenar la información, es un pilar para mejorar la experiencia de usuario y, por lo tanto, la de los jugadores en las campañas y sesiones.

No es una plataforma pensada para el juego remoto. Tampoco busca ser una simple herramienta de anotaciones para los jugadores, un sector ya muy explotado.

\subsubsection{¿Es un producto o servicio durable o estacional?}
Es un producto durable. Está concebido como una herramienta organizadora de anotaciones y registros de juego, que evoluciona con el tiempo mediante actualizaciones y la incorporación de nuevos manuales de juego. La oferta de la empresa para la comunidad de creadores de contenido también puede actualizarse según las necesidades y las opciones existentes.

\subsubsection{Características que influyen en la decisión de compra}
La característica más relevante para la decisión de compra es la facilidad de uso. El software busca simplificar el registro de información, en particular las anotaciones durante una sesión, una necesidad frecuente en las partidas, cuya falta hace perder el ritmo y compromete la inmersión y la experiencia de los jugadores.

También será relevante la comunidad: un buen catálogo de contenidos, un proceso ágil de publicación de contenidos nuevos y una remuneración correcta de los creadores permiten que la base de usuarios crezca por el solo hecho de que los usuarios sean proveedores de contenido además de consumir el de otros.

\subsubsection{¿Existen estudios que respalden las hipótesis?}
Se está realizando un estudio mediante encuestas, clasificadas según la experiencia en el juego para distinguir el tipo de usuario del software: jugadores nuevos, jugadores con experiencia y \textit{Dungeon Masters}. Su finalidad es corroborar las hipótesis sobre las experiencias en el rubro de los juegos de mesa, sobre las necesidades desatendidas y sobre la posibilidad de utilizar esta herramienta. Los resultados se incorporarán al plan cuando estén disponibles y se contrastarán, además, con los datos de uso real durante la validación con usuarios de la etapa de lanzamiento.

\subsubsection{Plan de lanzamiento}
El lanzamiento es escalonado y sigue las etapas del plan:
\begin{itemize}
  \item \textbf{Lanzamiento (primer año):} la aplicación web para DMs se presenta en Rosario en el mes 6, con el plan premium y los cobros habilitados; la aplicación móvil para jugadores se incorpora en el mes 9. Se prioriza la validación con usuarios reales.
  \item \textbf{Consolidación (segundo año):} se amplía la cobertura a la región, se refina el producto según la validación y se prepara el \textit{marketplace}.
  \item \textbf{Expansión (tercer año):} se lanza el \textit{marketplace} de contenido de la comunidad y se extiende la cobertura a todo el país.
\end{itemize}

\subsubsection{Ciclo de vida del producto}
Se prevé que el producto atraviese las siguientes fases durante el horizonte del plan:
\begin{itemize}
  \item \textbf{Introducción (etapa de lanzamiento):} baja cantidad de usuarios, con foco en dar a conocer el producto, validar la propuesta de valor y captar a los primeros DMs adoptantes.
  \item \textbf{Crecimiento (etapas de consolidación y expansión):} aumento sostenido de usuarios y ampliación de la oferta con la Línea Esencial y el \textit{marketplace}.
\end{itemize}
La madurez del producto queda fuera del horizonte del plan. Como líneas de evolución posteriores se consideran la incorporación de otros sistemas de juego y de integraciones con contenidos de terceros.

\subsubsection{Marca}
El producto se comercializa bajo la marca \nombreproyecto{}. Se definirá una identidad visual propia, con logotipo y paleta de colores, alineada con el posicionamiento descrito en la sección de metas de posicionamiento: cercana a la comunidad hispanohablante, accesible y orientada a la creatividad.

\subsection{Precio: estrategia de precio}

\subsubsection{Condicionantes del precio}
\begin{itemize}
  \item \textbf{Factores socioeconómicos locales:} el mercado inicial es Argentina, donde la inflación y las fluctuaciones del tipo de cambio afectan la capacidad de gasto de los usuarios en entretenimiento. Por eso los precios se definen en dólares estadounidenses, para preservar su valor, y pueden cobrarse en pesos argentinos según el tipo de cambio vigente, con una política regionalizada y accesible.
  \item \textbf{Disposición a pagar del segmento:} existe una baja disposición inicial a pagar por herramientas digitales dentro del pasatiempo, ya que parte de la comunidad está habituada a soluciones gratuitas o a métodos tradicionales (papel, mensajería).
  \item \textbf{Factores técnicos y operativos:} la estructura de precios debe cubrir los costos de alojamiento de datos, mantenimiento de la infraestructura y servicios externos de inteligencia artificial utilizados para asistir al DM.
  \item \textbf{Dinámica de red del producto:} como el DM atrae de cuatro a seis jugadores por mesa, cualquier barrera económica de entrada frenaría la adopción.
\end{itemize}

\subsubsection{Estrategia de precio}
La estrategia principal es un modelo \textit{freemium} B2C complementado con un esquema C2C (\textit{marketplace}).
\begin{itemize}
  \item \textbf{Plan gratuito (núcleo).} Acceso sin costo a las funcionalidades centrales para jugar: creación de campañas, hojas de personaje, notas y sincronización en tiempo real mediante WebSockets. Esto asegura una barrera de entrada nula para los DMs novatos y los ingresantes al pasatiempo.
  \item \textbf{Suscripción premium (SaaS B2C).} Orientada a los DMs experimentados. Ofrece características avanzadas, como un mayor límite de almacenamiento para mapas, un uso extendido del servicio de inteligencia artificial para diálogos y opciones avanzadas de personalización. Es una suscripción anual de pago único, con un precio en dólares que se escalona de forma progresiva a lo largo de las tres etapas del plan.
  \item \textbf{Comisiones del \textit{marketplace} (C2C).} A partir de su lanzamiento, en la tercera etapa, la plataforma retiene un porcentaje (por ejemplo, entre el 15\,\% y el 20\,\%) de las transacciones generadas por la venta de campañas, mapas y recursos creados por la comunidad.
\end{itemize}

\paragraph{Políticas de descuentos, bonificaciones y créditos.}
Se propone la siguiente política estandarizada, sujeta a validación durante la etapa de lanzamiento:
\begin{itemize}
  \item \textbf{Sin créditos ni cuotas:} la suscripción premium se abona al contado mediante pasarelas de pago digitales.
  \item \textbf{Bonificación por difusión:} se otorgan suscripciones premium sin costo, en cantidad limitada, a creadores de contenido y canales de transmisión que muestren la herramienta en uso real (canje por mención).
  \item \textbf{Promociones puntuales:} se admiten descuentos acotados en eventos de lanzamiento y convenciones, con un tope y una vigencia definidos de antemano.
\end{itemize}

\subsection{Comunicaciones integradas de marketing}

\subsubsection{Planificación estratégica de las acciones de comunicación}
\paragraph{Acción 1: presencia en eventos locales (fase 1, Rosario).}
\begin{itemize}
  \item \textbf{Objetivo:} dar a conocer el MVP en el mercado local, validar funcionalidades y captar a los primeros DMs.
  \item \textbf{Tipo:} promoción directa y relaciones públicas.
  \item \textbf{Estilo:} demostrativo e interactivo.
  \item \textbf{Medio:} \textit{stands} físicos y demostraciones en vivo en convenciones como Rosario Juega Rol y Crack Bang Boom.
  \item \textbf{Costo estimado:} bajo a moderado (costo del \textit{stand}, folletería y \textit{merchandising} básico).
\end{itemize}

\paragraph{Acción 2: patrocinio de \textit{campaign playthroughs} (fases 2 y 3, regional y nacional).}
\begin{itemize}
  \item \textbf{Objetivo:} persuadir e informar a una audiencia masiva de jugadores mostrando la herramienta en uso real.
  \item \textbf{Tipo:} publicidad digital y \textit{marketing} de influencia.
  \item \textbf{Estilo:} inmersivo y testimonial (el DM del canal de transmisión utiliza la aplicación en vivo).
  \item \textbf{Medio:} canales de Twitch y YouTube dedicados a TTRPG en Argentina.
  \item \textbf{Costo estimado:} moderado (canje de suscripciones premium sin costo o pagos por mención publicitaria).
\end{itemize}

\subsubsection{Diseño del sitio web (7 C y pantallas descriptivas)}
El diseño del sitio web se organiza según las siete C del diseño de sitios:
\begin{itemize}
  \item \textbf{Contexto:} una interfaz inmersiva pero limpia, que evita la sobrecarga de información que enlentece la experiencia de juego.
  \item \textbf{Contenido:} manuales de reglas base (D\&D 5e), material oficial de la Línea Esencial y contenido generado por los usuarios.
  \item \textbf{Comunidad:} un sistema de foros o valoraciones en el \textit{marketplace}, donde los usuarios reseñan el contenido creado por otros (C2C).
  \item \textbf{Customización:} capacidad del DM de personalizar las reglas y adaptar el RuleSet base.
  \item \textbf{Comunicación:} sistema de notificaciones e invitaciones directas para que los jugadores se unan a las campañas.
  \item \textbf{Conexión:} interacción en tiempo real (WebSockets) que refleja al instante el daño o el uso de objetos durante la partida presencial.
  \item \textbf{Comercio:} un entorno integrado del \textit{marketplace} para transacciones rápidas y seguras de módulos narrativos.
\end{itemize}

\subsection{Distribución}

\subsubsection{Factores condicionantes de la distribución y principales canales}
Al ser un producto puramente digital (\textit{Software as a Service}), la distribución no está condicionada por la logística física, el transporte o el almacenamiento de bienes. Los condicionantes principales son la disponibilidad y la latencia de los servidores, y las políticas de las tiendas de aplicaciones y de los motores de búsqueda (SEO), que actúan como intermediarios de la visibilidad.

\subsubsection{Estrategia de distribución}
La estrategia es de distribución directa y selectiva en su fase inicial. El producto se entrega directamente desde los servidores de la empresa al dispositivo del usuario, mediante la web o la aplicación móvil, sin intermediarios físicos.

\subsubsection{Canales (tipo y nivel)}
\begin{itemize}
  \item \textbf{Tipo:} canal digital directo.
  \item \textbf{Nivel:} nivel 0 (productor a consumidor final). No existen mayoristas ni minoristas tradicionales. El \textit{marketplace} integrado, una vez lanzado, también opera un canal directo de nivel 0 entre los creadores de contenido y los DMs y jugadores.
\end{itemize}

\subsubsection{Análisis de localización de puntos de venta propios}
Al tratarse de un producto de base tecnológica y distribución digital, el concepto tradicional de punto de venta se adapta a un modelo híbrido, que combina un canal digital centralizado con activaciones físicas estratégicas.

\paragraph{Punto de venta digital (principal).}
La comercialización, distribución y operación del software se centralizan en una página web propia, alojada inicialmente en infraestructura propia, con escalado a la nube según la cantidad de usuarios. Este dominio actúa como la sucursal virtual y punto de venta oficial de \nombreproyecto{}: en él los usuarios (DMs y jugadores) pueden registrarse, acceder a las funcionalidades del plan gratuito, contratar las suscripciones premium y, una vez lanzado el \textit{marketplace}, realizar transacciones. La aplicación móvil se distribuye, además, a través de las tiendas de aplicaciones. Su localización es de alcance global y disponibilidad permanente, y se apoya en una fuerte optimización para motores de búsqueda (SEO) para captar el tráfico de usuarios hispanohablantes interesados en herramientas de TTRPG.

\paragraph{Puntos de venta físicos efímeros (eventos locales).}
Para traccionar la adopción inicial y superar la barrera de desconfianza tecnológica, la empresa instala puntos de venta físicos temporales: \textit{stands} y participación activa en eventos locales, clubes de rol y convenciones de cultura popular (como Rosario Juega Rol y Crack Bang Boom, entre otros). Estos \textit{stands} funcionan como terminales de demostración y adquisición directa: permiten a los DMs probar la herramienta en vivo durante partidas de exhibición y crear su cuenta en el acto escaneando códigos QR. Esta presencia física garantiza el contacto directo con el nicho y acelera la validación del producto y la captación de los primeros usuarios de la región.

\subsubsection{Opciones de asociación}
\begin{itemize}
  \item \textbf{Clubes de rol y tiendas de juegos de mesa locales:} especialmente en Rosario y la Región Centro, con códigos QR en las mesas físicas que dirijan a la descarga de la aplicación.
  \item \textbf{Creadores de contenido independientes:} ilustradores y escritores de historias que puedan subir su material al \textit{marketplace} de forma exclusiva durante su lanzamiento, en la tercera etapa del plan.
\end{itemize}

\subsubsection{Gestión del JIT (\textit{just in time})}
Como \nombreproyecto{} es un producto de base tecnológica, el concepto de \textit{just in time} se aplica a la arquitectura del software y a la gestión de los recursos:
\begin{itemize}
  \item \textbf{Subsistema de producción (desarrollo):} metodologías ágiles y flujos de integración y despliegue continuo (CI/CD) permiten lanzar actualizaciones, corregir errores y desplegar nuevas funcionalidades con rapidez, según la retroalimentación validada de los usuarios.
  \item \textbf{Subsistema de almacenamiento y distribución:} la infraestructura inicial es propia (\textit{on-premise}) y, a medida que crezca la cantidad de usuarios, la nube permite un escalado dinámico. No se fabrican licencias por adelantado: el servidor asigna recursos (ancho de banda, almacenamiento de imágenes y bases de datos) bajo demanda y en tiempo real a medida que nuevos grupos inician sus sesiones de juego.
\end{itemize}
